\documentclass[11pt,a4paper]{article}
\usepackage[margin=20mm]{geometry}
\usepackage{amsmath,amssymb,mathtools}
\usepackage{enumitem}
\usepackage{xcolor}
\usepackage{fontspec}
\setmainfont{Arial Unicode MS}
\setmonofont{Menlo}
\definecolor{navy}{HTML}{172554}
\definecolor{blue}{HTML}{1D4ED8}
\setlength{\parskip}{0.65em}
\setlength{\parindent}{0pt}
\newcommand{\weektitle}[1]{\section*{\color{blue}#1}}
\newcommand{\question}[1]{\subsection*{#1}}
\newcommand{\answerbox}[1]{\vspace*{#1}}
\newcommand{\code}[1]{\texttt{#1}}
\begin{document}

\begin{center}
{\Large\color{navy} Transformer From Scratch Lab}\par
\vspace{9mm}
{\huge\color{navy} Week 0 Problem Set}\par
\vspace{4mm}
{\Large 先看懂信息怎样流动}
\end{center}

\weektitle{课程位置}
本题集不要求实现 attention。目标是让你在 Week 1 前，能准确追踪 token 表、矩阵乘法、点积和 stable softmax 的输入、输出与 shape。

\textbf{提交物：}书面解答、Shape Checker Lab、一个真实错误或主动注入错误的诊断记录，以及闭卷信息流图。

\textbf{使用规则：}Part A 与 Part B 必须在进入 Lab 前独立完成。Part C 必须在 Lab 后引用自己的输入、输出、测试或评分反馈。

\newpage
\weektitle{Part A：概念与推导}

\question{A1. Token 表与 shape（推导）}
令 $X \in \mathbb{R}^{3 \times 2}$ 表示三个 token 的二维表示，令 $W \in \mathbb{R}^{2 \times 4}$。
\begin{enumerate}[label=(\arabic*)]
\item 写出 $XW$ 的 shape。
\item 用一句话说明输出的每一行和每一列分别代表什么。
\item 令 $X_{2,1}$ 增加 $1$。输出矩阵的哪些元素可能变化？为什么？
\end{enumerate}
\answerbox{34mm}

\question{A2. 矩阵乘法不是逐元素乘法（推导）}
给定
\[
A = \begin{bmatrix} 1 & 2 \\ 0 & 3 \\ 4 & 1 \end{bmatrix},
\qquad
B = \begin{bmatrix} 2 & 1 \\ 1 & 0 \end{bmatrix}.
\]
\begin{enumerate}[label=(\arabic*)]
\item 手算 $AB$ 的第 $1$ 行与第 $3$ 行。
\item 写出一般元素 $(AB)_{i,j}$ 的计算规则。
\item 解释为什么这个规则允许每个 token 经过同一组权重后获得新的特征。
\end{enumerate}
\answerbox{34mm}

\question{A3. Dot product 与匹配分数（推导）}
令 $u = [1,2]$，$v = [3,-1]$，$w = [1,1]$。
\begin{enumerate}[label=(\arabic*)]
\item 计算 $u \cdot v$ 与 $u \cdot w$。
\item 若 $u$ 是 query，$v$ 和 $w$ 是两个 key，哪个 key 在这个极小例子中更匹配？
\item 解释为何这个分数还不能直接作为读取权重。
\end{enumerate}

\newpage
\question{A4. Stable softmax（推导）}
对 $z=[2,1,0]$：
\begin{enumerate}[label=(\arabic*)]
\item 写出 $\operatorname{softmax}(z)_i = \dfrac{e^{z_i}}{\sum_j e^{z_j}}$ 的三个分量，并验证它们的和为 $1$。
\item 证明对任意常数 $c$，$\operatorname{softmax}(z)=\operatorname{softmax}(z+c\mathbf{1})$。
\item 用第 $(2)$ 问解释稳定实现为何可以先计算 $z-\max_i z_i$。
\end{enumerate}
\answerbox{45mm}

\weektitle{Part B：使用边界与全局位置}
\question{B1. 不合法的 shape（反例）}
构造两个二维矩阵 $A \in \mathbb{R}^{m \times n}$ 与 $B \in \mathbb{R}^{p \times q}$，使得 $n \ne p$。
\begin{enumerate}[label=(\arabic*)]
\item 明确写出左右内维度。
\item 写出 Shape Checker 应报告的错误类别和错误信息必须包含的内容。
\item 解释若只让 NumPy 抛出原始异常，学习者会缺少哪一条机制信息。
\end{enumerate}
\answerbox{35mm}

\question{B2. Softmax 的错误轴（反例）}
设 score matrix 为 $S \in \mathbb{R}^{3 \times 3}$，其中第 $i$ 行表示 query $i$ 对三个 key 的分数。
\begin{enumerate}[label=(\arabic*)]
\item 解释沿最后一维归一化意味着谁在竞争权重。
\item 解释沿第 $0$ 维归一化时谁在竞争权重。
\item 构造一个最小矩阵 $S$，使两种归一化产生显著不同的解释。
\end{enumerate}

\newpage
\question{B3. Week 0 在整体中的位置（全局连接）}
画出下列信息流，并为每个箭头标出 Week 0 已掌握的操作：
\[
\text{token table}
\rightarrow Q,K,V
\rightarrow QK^\top
\rightarrow \operatorname{softmax}\!\left(\frac{QK^\top}{\sqrt{d_k}}\right)
\rightarrow \operatorname{weighted\ values}.
\]
最后用 $150$--$200$ 字解释：为什么 Week 0 还没有真正实现 attention，却已为 Week 1 的每一步建立前置能力？
\answerbox{52mm}

\weektitle{Part C：Lab 后工程问题}
\question{C1. Shape Checker 的契约（工程设计）}
为什么要将 shape 验证与输出 shape 推导拆成两个函数？\\
分别为 \code{validate\_matrix\_shape} 和 \code{matrix\_product\_shape} 写出一个它单独负责的错误场景。
\answerbox{25mm}

\question{C2. 一个真实错误（工程诊断）}
选择一次 Lab 失败或主动注入错误 shape。\\
记录输入、现象、错误类别、原因假设、最小验证测试，以及修复后理解的变化。
\answerbox{35mm}

\question{C3. 从 Lab 到 Week 1（工程迁移）}
说明 Shape Checker 如何预先防止 $Q=XW_Q$、$K=XW_K$、$V=XW_V$ 三次投影的错误；再说明它无法防止哪一种 shape 正确但机制仍错误的问题。

\newpage
\weektitle{闭卷自检与提交清单}
不看材料，用 $3$--$5$ 分钟解释：
\begin{enumerate}[label=(\arabic*)]
\item 一个 $(3,2)$ token 表经过 $(2,4)$ 投影后发生了什么？
\item dot product、softmax 与加权和在 Week 1 attention 中分别做什么？
\item 给出一种 shape 正确但机制错误的 attention 实现方式。
\end{enumerate}

\vspace{8mm}
\begin{tabular}{ll}
$\square$ & Part A 与 Part B 已在 Lab 前完成。\\[3mm]
$\square$ & Shape Checker 的公开与 hidden tests 已通过。\\[3mm]
$\square$ & Part C 引用了自己的 Lab 证据。\\[3mm]
$\square$ & 已记录一个真实错误或主动注入错误的诊断过程。\\[3mm]
$\square$ & 能闭卷画出 Week 0 到 Week 1 的信息流。\\
\end{tabular}

\vfill
\begin{center}\small Week 0 Problem Set\end{center}
\end{document}
